\begin{abstract}
Using the written statements of OpenAI's \emph{The Quasi-Riemann Hypothesis: A Zero-Free Half-Plane $\Rea s>7/8$} (30 September 2026) as input, we refine its compensated-probe argument. Every Dirichlet $L$-function, every finite-order Hecke $L$-function over $\Q(\sqrt{-3})$, and the Riemann zeta function have no zeros in $\Rea s>\sigma_\dagger$, where $\sigma_\dagger=0.8749570194\ldots$ is the root in $(0.8749570194,0.8749570195)$ of $7884s^3-18819s^2+14643s-3686$. In particular, $\Rea s>43747851/50000000$ is zero-free; principal poles at $s=1$ are allowed. Optimising the skew parameter gives the quadratic endpoint $(1507-2\sqrt{921})/1653$. Extending the plain fourth-moment range from $[3/4,1]$ to $[37/50,1]$ gives the cubic endpoint $\sigma_\dagger$. Exact algebraic certificates are included as ancillary files. One mean-value hypothesis, stated in Section~\ref{sec:limits} and not proved here, would give the boundary $328/375$.
\end{abstract}

\section{Introduction}
OpenAI's paper \cite{openai2026q} proves a common zero-free half-plane for the finite-order Hecke family over the Eisenstein field and transfers it to all Dirichlet $L$-functions. Its completed cubic reflection, compensated probe, marked inverse moment, and plain fourth moment provide the analytic structure used here. We retain that structure and enlarge its geometric and moment parameters.

Taking the written statements of \cite{openai2026q} as input, with the hypotheses specified below, we prove:

\begin{theorem}\label{thm:main}
Let $\sigma_\dagger$ be the unique root in
\[
 \frac{8749570194}{10^{10}}<\sigma_\dagger<\frac{8749570195}{10^{10}}
\]
of
\begin{equation}\label{eq:G}
 G(s)=7884s^3-18819s^2+14643s-3686.
\end{equation}
Every finite-order Hecke $L$-function over $\F=\Q(\sqrt{-3})$ and every Dirichlet $L$-function, of every modulus, have no zeros in $\Rea s>\sigma_\dagger$. This includes $\zeta(s)$. The pole of a principal character at $s=1$ is allowed.
\end{theorem}

\begin{corollary}\label{cor:main}
The same functions have no zeros in
\[
 \Rea s>\frac{43747851}{50000000}
 =\frac78-\frac{2149}{50000000}.
\]
\end{corollary}

The boundary $20999/24000=7/8-1/24000$ of Theorem~\ref{thm:rational2} was obtained independently by Z.~Gao in a draft dated 7 October 2026 \cite{gao2026}, with the same parameters $\ell=1001/6000$ and $b=1/8$. Gao also identifies the fixed-skew envelope's limiting boundary as approximately $0.8749572006$. Here we optimise the skew parameter $b$, obtaining the endpoint $\sigma_*=(1507-2\sqrt{921})/1653$ in Theorem~\ref{thm:quadratic}, and extend the plain fourth-moment range to $37/50\le\kappa\le1$, obtaining $\sigma_\dagger$ in Theorem~\ref{thm:cubic}.

The input is OpenAI's paper \cite{openai2026q}, dated 30 September 2026; references to that paper use repository \nolinkurl{openai/math} at commit \nolinkurl{adc7f1241b42e322a6451854ab7e4b4c146bf78a}. Its final theorems---the $7/8$ half-plane for $\zeta$, all Dirichlet $L$-functions and all finite-order Hecke $L$-functions over $\Q(\sqrt{-3})$---are covered by the Lean formalization described in \cite{openai2026lean}, with principal poles excluded. The intermediate statements of \cite{openai2026q} used here are taken as written, with their hypotheses listed in Appendix~\ref{app:inputs}. In particular, its Theorem~1.1 supplies $\beta^*\le7/8$. Part~II of \cite{openai2026q} is written under the standing supposition $\beta^*>7/8$; Appendix~\ref{app:partii} proves that the auxiliary results used here do not depend on that supposition, retaining their local hypotheses. We change the geometry of the compensated probe and re-derive the completed norm and low estimate, Euler-factor domain, row count, contours and heights that depend on it. Theorem~\ref{thm:fourth} extends the plain fourth-moment estimate of \src{Lemma 18.1}, and the subsequent parameter optimisation gives an improvement of $4.298\ldots\times10^{-5}$ over $7/8$ within OpenAI's method.

Each stage is a separate theorem: the geometric criterion (Theorem~\ref{thm:geometry}), the general row count (Theorem~\ref{thm:count}), the rational stages (Theorems~\ref{thm:rational1} and \ref{thm:rational2}), the quadratic boundary (Theorem~\ref{thm:quadratic}), the extended fourth moment (Theorem~\ref{thm:fourth}), and the cubic boundary (Theorem~\ref{thm:cubic}). Section~\ref{sec:limits} states the limit of the specified envelope model and states one unproved mean-value hypothesis under which the method would give $328/375$.

All half-planes in this paper are strict. Their boundaries are common to the entire character family. Constants, smooth orders, excluded prime sets, and lower thresholds for a fixed target may depend on that target.
